\documentclass[12pt]{article}
\usepackage[T1]{fontenc}
\usepackage[utf8]{inputenc}
\usepackage{lmodern}
\usepackage{textcomp}
\usepackage{enumitem}
\usepackage{geometry}
\geometry{margin=1in}
\begin{document}
\begin{center}
Answer Key - Political Science - Graduate Level Assessment 15 \
\small (Answers are ordered exactly as in the question paper.)
\end{center}

\section*{Multiple Choice}
\begin{enumerate}[label=\arabic*.]
\item \textbf{Answer:} D
\\ \textit{Options:} A) Leaders have incentives to lie; B) If leaders revealed their programs, they would be more likely to be attacked; C) Leaders will not always grant foreign monitors access to their nuclear programs; D) ALL of the above
\\ \textit{Note:} The uncertainty surrounding which states possess nuclear weapons stems from a combination of factors, including varying degrees of transparency among nations, differing national security policies, and the complexities of verification and intelligence gathering, all of which contribute to the ambiguity in global nuclear arsenals.
\item \textbf{Answer:} B
\\ \textit{Options:} A) Sovereignty.; B) Identity.; C) All of these options.; D) Hegemonic ethnicity.
\\ \textit{Note:} Identity is essential to societal security because it shapes group cohesion, social stability, and the collective sense of belonging, which are crucial for maintaining peace and mitigating conflict within diverse populations.
\item \textbf{Answer:} B
\\ \textit{Options:} A) The US inhibited the marketization of the Russian economy; B) The US promoted the marketization of the Russian economy; C) The US supported public ownership of natural resources; D) None, the US was only concerned with security issues
\\ \textit{Note:} The correct answer is supported by the fact that the Clinton Administration actively advocated for and implemented policies aimed at transitioning Russia from a state-controlled economy to a market-based system, which included promoting privatization, deregulation, and the establishment of free market practices.
\item \textbf{Answer:} A
\\ \textit{Options:} A) 1970 s; B) 1990 s; C) 1960 s; D) 1980 s
\\ \textit{Note:} The cyber-security discourse began to emerge in the 1970 s as it was during this decade that the proliferation of computer technology and the recognition of vulnerabilities in digital systems prompted serious discussions among politicians, academics, and industry experts about the need for security measures and policies to protect sensitive information and infrastructure.
\item \textbf{Answer:} C
\\ \textit{Options:} A) The arms trade shifts from being a private to a government controlled enterprise.; B) An increase in the defence trade.; C) A decrease in the arms trade.; D) A growth in the number and variety of weapons traded.
\\ \textit{Note:} The correct answer is that a decrease in the arms trade did not emerge as a trend in the 20 th century because, in fact, the arms trade significantly increased due to factors such as the Cold War, technological advancements, and the rise of military industrial complexes globally.
\end{enumerate}

\section*{True / False}
\begin{enumerate}[label=\arabic*.]
\setcounter{enumi}{5}
\item \textbf{Answer:} True
\\ \textit{Options:} A) True; B) False
\\ \textit{Note:} The statement is true because ideational approaches prioritize the influence of ideas, beliefs, and ideological frameworks-such as democracy and communism-over material factors like economic or military interests in shaping U.S. foreign policy during the Cold War.
\item \textbf{Answer:} False
\\ \textit{Options:} A) True; B) False
\\ \textit{Note:} Critics actually contend that the human security concept is too broad and encompasses a wide range of issues, thereby addressing global challenges comprehensively rather than narrowly.
\item \textbf{Answer:} False
\\ \textit{Options:} A) True; B) False
\\ \textit{Note:} The statement is false because NATO had previously employed military force during the Bosnian War in the mid-1990 s, specifically through airstrikes and operations aimed at enforcing peace agreements.
\item \textbf{Answer:} True
\\ \textit{Options:} A) True; B) False
\\ \textit{Note:} The statement is true because the Security Council holds significant authority to make binding decisions regarding international peace and security, including the power to impose sanctions and authorize military action, which positions it as the most powerful body within the United Nations framework.
\item \textbf{Answer:} False
\\ \textit{Options:} A) True; B) False
\\ \textit{Note:} Natural resource trade often leads to authoritarianism and the "resource curse," where wealth from resources concentrates power in the hands of elites, undermining political pluralism and civil liberties.
\end{enumerate}

\section*{Long Form Response}
\begin{enumerate}[label=\arabic*.]
\setcounter{enumi}{10}
\item \textbf{Answer:} The political subjecthood of child soldiers is a critical lens through which to understand their role in conflict situations, as it recognizes them as active participants rather than mere victims. This perspective challenges the dominant narrative that portrays children solely as innocent bystanders, highlighting their agency and the complex realities they navigate within warfare. By acknowledging their political subjecthood, we can better address the socio-political factors that lead to their recruitment and the implications of their involvement in armed conflict. This shift in perception invites a reevaluation of policies and interventions that often overlook the voices and experiences of child soldiers, thereby fostering a more nuanced understanding of their plight and potential for advocacy. Ultimately, recognizing child soldiers as political subjects enriches discussions on accountability, rehabilitation, and the broader impact of war on youth.
\\ \textit{Note:} correct because it emphasizes the recognition of child soldiers as active agents in conflict, thereby challenging the simplistic victim narrative and prompting a deeper examination of the socio-political dynamics that contribute to their involvement in warfare.
\item \textbf{Answer:} The Great Law of Peace of the Haudenosaunee Confederacy is significant for its foundational principles of democracy, consensus-building, and collective governance, which have profoundly influenced both indigenous and modern political thought. It emphasizes the importance of unity among diverse nations, promoting peace and mutual respect while establishing a framework for decision-making that prioritizes the welfare of the community. The impact of these principles is evident in their ability to foster social cohesion and stability, serving as a model for governance that values dialogue over conflict. Additionally, the Great Law of Peace has contributed to discussions on indigenous rights and sovereignty, highlighting the relevance of traditional governance systems in contemporary political discourse. Overall, its legacy underscores the importance of inclusive governance and the potential for indigenous frameworks to inform broader political practices.
\\ \textit{Note:} This answer is correct because it articulates how the Great Law of Peace embodies democratic principles and collective governance that foster unity and community welfare, illustrating its significant influence on both indigenous and contemporary political frameworks.
\end{enumerate}

\vfill
\noindent\textit{End of Answer Key}
\end{document}